% Scope: Derive reception and the MRI forward model, then identify exactly when Fourier inversion is justified; establish PSF and complex-noise language.
% Status: architecture only; no chapter prose or completed problems yet.
% Prerequisites: Chapters 1 through 5.
% Planned worked examples: Two-point object, finite readout blur and noise through a DFT.
% Planned figures: Receive sensitivity, Fourier pairs and point-spread functions.
% Planned challenge problems: Derive the forward operator for a decaying readout; compare sampling and resolution.
\chapter{Signal Formation and Fourier Imaging}
\label{ch:06}

\section{From precession to received voltage}
\label{sec:06:from-precession-to-received-voltage}
\subsection{Faraday induction and reciprocity}
\subsection{Complex demodulation and receiver phase}
\subsection{Receive sensitivity and signal normalization}

\section{The MRI signal equation}
\label{sec:06:the-mri-signal-equation}
\subsection{Spatial integration of transverse magnetization}
\subsection{Encoding, relaxation and off-resonance during readout}
\subsection{Motion, multiple species and model limitations}

\section{The Fourier imaging limit}
\label{sec:06:the-fourier-imaging-limit}
\subsection{Stationary object and time-independent weighting}
\subsection{Forward and inverse transform conventions}
\subsection{Discrete sampling and voxel integration}

\section{Image formation and point-spread functions}
\label{sec:06:image-formation-and-point-spread-functions}
\subsection{Finite support and truncation}
\subsection{Apodization, convolution and effective resolution}
\subsection{Center and periphery without oversimplification}

\section{Incomplete and modified Fourier data}
\label{sec:06:incomplete-and-modified-fourier-data}
\subsection{Partial Fourier and conjugate-symmetry assumptions}
\subsection{Phase estimation and homodyne concepts}
\subsection{Zero padding, interpolation and display grids}

\section{Complex images and noise}
\label{sec:06:complex-images-and-noise}
\subsection{DFT normalization and covariance}
\subsection{Phase images, magnitude images and noise floors}
\subsection{Receive weighting and the meaning of image intensity}

\section{Worked examples}
\label{sec:06:worked-examples}
% Planned calculations: Two-point object, finite readout blur and noise through a DFT.
\subsection{Derivation and quantitative calculation}
\subsection{Assumptions, limiting cases and scanner interpretation}

\section{Challenge problems}
\label{sec:06:challenge-problems}
% Problem design: Derive the forward operator for a decaying readout; compare sampling and resolution.
% Use challengeproblem and stable labels prob:06:<topic>.
% Complete solutions belong in Appendix E, section for this chapter.
% Use \begin{solution}{prob:06:<topic>} ... \end{solution} there.
\subsection{Derivation and quantitative design}
\subsection{Model criticism and integrated reasoning}
